\documentclass[10pt,letterpaper]{article}
\usepackage[margin=0.45in]{geometry}
\usepackage[hidelinks]{hyperref}
\usepackage{enumitem}
\usepackage{titlesec}
\usepackage{tabularx}
\usepackage{array}
\usepackage[T1]{fontenc}
\usepackage[utf8]{inputenc}
\usepackage{lmodern}
\usepackage{setspace}
\setstretch{1.07}
\pagestyle{empty}
\setlength{\parindent}{0pt}
\setlength{\tabcolsep}{0pt}
\titleformat{\section}{\large\bfseries}{}{0em}{}[\titlerule]
\titlespacing*{\section}{0pt}{6pt}{4pt}
\newcommand{\resumeItem}[1]{\item #1}
\newcommand{\resumeExpHeading}[3]{% 
  \begin{tabularx}{\textwidth}{X r}
    \textbf{#1} $\mid$ \textit{#2} & #3 \\
  \end{tabularx}\vspace{0pt}
}
\newcommand{\resumeEduHeading}[3]{% 
  \begin{tabularx}{\textwidth}{X r}
    \textbf{#1} $\mid$ \textit{#2} & #3 \\
  \end{tabularx}\vspace{0pt}
}
\newcommand{\resumeProjectHeading}[3]{% 
  \begin{tabularx}{\textwidth}{X r}
    \textbf{#1} \textit{#2} & #3 \\
  \end{tabularx}\vspace{0pt}
}
\newenvironment{resumeItemList}{\begin{small}\begin{itemize}[leftmargin=*, itemsep=2pt, topsep=3pt, parsep=0pt]}{\end{itemize}\end{small}\vspace{4pt}}
\begin{document}
\begin{tabularx}{\textwidth}{@{}X r@{}}
    {\fontsize{18}{28}\selectfont \textbf{Advait Paranjpe}} & \href{mailto:advait_paranjpe@outlook.com}{Email: advait\_paranjpe@outlook.com} \\[-1pt]
    \href{https://github.com/AdvaitParanjpe}{GitHub: github.com/AdvaitParanjpe} & Mobile: +1 (213) 284-9936 \\[-1pt]
    \href{https://linkedin.com/in/advait-paranjpe}{LinkedIn: linkedin.com/in/advait-paranjpe} & \href{https://advaitparanjpe.com}{Website: advaitparanjpe.com}
\end{tabularx}
\vspace{3pt}
\section*{Education}
\resumeEduHeading{University of California, Los Angeles}{M.S. in Electrical and Computer Engineering}{Sep 2025 -- Dec 2026}
\begin{resumeItemList}
    \resumeItem{Coursework: CPU/GPU Architecture, Memory Systems, Cache Coherence, Interconnect Design.}
    \resumeItem{GPA: 3.7/4.0}
\end{resumeItemList}
\resumeEduHeading{The University of Manchester}{B.Eng. in Electrical and Electronic Engineering}{Sep 2022 -- Jul 2025}
\begin{resumeItemList}
    \resumeItem{Dissertation: ``Cascaded PLL Microwave Synthesizer for 5G systems.''}
    \resumeItem{\textbf{First Class Honours} -- highest UK degree classification.}
\end{resumeItemList}
\section*{Experience}
\resumeExpHeading{Jaguar Land Rover North America}{Incoming Semiconductor Applications Engineering Intern}{Jul 2026 -- Sep 2026}
\resumeExpHeading{VAST Lab, UCLA}{Student Researcher -- GPU/FPGA Acceleration}{Jan 2026 -- Mar 2026}
\begin{resumeItemList}
    \resumeItem{Profiled ROCm attention kernels and memory bottlenecks to identify FPGA acceleration opportunities for Engram LLM inference.}
    \resumeItem{Modeled FPGA cache architectures, increasing projected hit rate from $\sim$5\% at 256 entries to $\sim$48\% with a 64k-entry, 4-way design.}
    \resumeItem{Translated workload profiling into FPGA tradeoffs across cache capacity, lookup throughput, locality, and on-chip memory usage.}
\end{resumeItemList}
\resumeExpHeading{NESO / National Grid ESO}{Software Engineering Intern (two summers)}{Jul 2024 -- Sep 2025}
\begin{resumeItemList}
    \resumeItem{Re-architected a 100k+ LOC simulation engine through vectorization and pipeline optimization, cutting runtime by 10x.}
    \resumeItem{Optimized 100M+ records through chunked, memory-efficient execution, cutting RAM from 32 GB to 16 GB and runtime by 11x.}
    \resumeItem{Reduced CPU context-switching overhead through parallel API batching and staged processing, increasing throughput by 24x for time-sensitive market telemetry.}
\end{resumeItemList}
\section*{Projects}
\resumeProjectHeading{tinyNPU}{$\mid$ Matrix-Multiply Accelerator with ASIC Flow (SystemVerilog)}{May 2026 -- Jun 2026}
\begin{resumeItemList}
    \resumeItem{Designed synthesizable RTL for a 4x4 signed-INT8 matrix-multiply accelerator with a parallel MAC array, scratchpad memories, DMA data movement, and AXI/APB wrapper variants.}
    \resumeItem{Built double-buffered memories and independent load, compute, and output FSMs, achieving 64 cycles/tile steady-state throughput.}
    \resumeItem{Improved 100 MHz WNS from -5.83 ns to -0.34 ns through STA-guided MAC pipelining, then ran an OpenLane/SKY130 RTL-to-GDS flow producing DRC/LVS-clean layouts.}
\end{resumeItemList}
\resumeProjectHeading{Sparrow-V}{$\mid$ Sparse-Aware RISC-V Processor (SystemVerilog, C, Python)}{Mar 2026 -- Jun 2026}
\begin{resumeItemList}
    \resumeItem{Designed a synthesizable RV32I processor with a tightly coupled four-lane INT8/INT16 vector engine for quantized edge workloads.}
    \resumeItem{Implemented a custom vector ISA, pipeline control, forwarding, branch redirection, precise traps, memory backpressure, and 2:4 structured-sparse execution.}
    \resumeItem{Evaluated microarchitectural behavior through differential testing, randomized workloads, assertions, cycle counters, and synthesis-oriented area/timing analysis.}
\end{resumeItemList}
\resumeProjectHeading{Sparrow-Cluster}{$\mid$ Four-Core Coherent RISC-V SoC (SystemVerilog)}{Jun 2026}
\begin{resumeItemList}
    \resumeItem{Architected a four-core RV32I system with private 2 KiB L1 instruction/data caches, shared memory, and blocking cache operation.}
    \resumeItem{Designed a snoopy MSI protocol with write-back, write-allocate L1 data caches and a round-robin shared interconnect.}
    \resumeItem{Implemented invalidation/writeback handling, multicore synchronization, LR/SC, and performance counters, then validated scaling and false-sharing-sensitive workloads across 1, 2, and 4 cores.}
\end{resumeItemList}
\vspace{-0.5em}
\section*{Skills}
\begin{resumeItemList}
    \resumeItem{\textbf{Languages:} SystemVerilog, Verilog, VHDL, Tcl, Python, C, C++, Bash, SQL, MATLAB.}
    \resumeItem{\textbf{Digital Design:} Processor microarchitecture, RTL design, pipelining, cache coherence, interconnects, synthesis, STA, timing closure, and PPA analysis.}
    \resumeItem{\textbf{Tools:} Verilator, Vivado, Yosys, OpenLane, GTKWave, Git, Linux, Docker.}
\end{resumeItemList}
\end{document}
